\color{uniblue}\section[Выбор уставок токовой отсечки трансформатора]{\sloppy Выбор уставок токовой отсечки \mbox{трансформатора}} 
\color{black}

\begin{enumerate}[label=\arabic{section}.\arabic*, labelsep=4pt, leftmargin=0pt, itemindent=57pt]

\item
По условию селективности ток срабатывания отсечки выбирается большим максимального значения тока при КЗ на стороне низшего напряжения защищаемого силового трансфоматора по формуле:

\begin{equation}
\label{eq:to1}
I_{\textsl{с.з.ТО}} = k_{\textsl{н}} \cdot I_{\textsl{к.макс.}}^{(3)} ,
\end{equation}
\begin{eqexpl}[15mm]
\item{$I_{\textsl{к.макс.}}^{(3)}$} максимальное значение тока трехфазного КЗ на стороне НН, определяется при максимальном режиме питающей системы (сопротивление системы минимальное возможное), для трансформаторов с регулированием напряжения при расчете необходимо также учитывать минимальное сопротивление трансформатора при крайнем положении его регулятора напряжения, А;
\item{$k_{\textsl{н}}$} коэффициент надёжности для токовых отсечек без выдержки времени, установленных на понижающих трансформаторах, при использовании цифровых реле, может быть принят равным 1,2.
% у Радиуса принимается 1.2.
%\\192.168.11.240\Company\Ivanovo\Документация ЮНИТ М300\Разработка\Схемы ФБ ЮНИТ-М3\Трансформатор\ИЭУ Т 35 кВ Россети\07. Инфо по уставкам\Vybor-ustavok-Sirius_UV.pdf#nameddest=ТО_трансформатора
\end{eqexpl}

\item
Кроме условия (\ref{eq:to1}), должна быть обеспечена отстройка токовой отсечки от бросков тока намагничивания силовых трансформаторов по выражению:
\begin{equation}
    \label{eq:to2}
    I_{\textsl{с.з.ТО}} = k_{\textsl{броска}} \cdot I_{\textsl{ном.тр.}},
\end{equation}
\begin{eqexpl}[15mm]
    \item {$k_{\textsl{броска}}$} коэффициент, учитывающий увеличение тока при броске тока намагничивания, может быть принят равным 5,0;
    \item {$I_{\textsl{ном.тр.}}$} номинальный ток трансформатора, А.
\end{eqexpl} 
Эти броски возникают в момент включения под напряжение ненагруженного трансформатора и могут первые несколько периодов превышать номинальный ток трансфоматора в 5-7 раз. При расчете токовой отсечки для трансформатора по условию (\ref{eq:to1}) одновременно обычно выполняется и отстройка от броска тока намагничивания этого трансформатора.

\item
Для оценки чувствительности токовой отсечки необходимо определить коэффициент чувствительности по формуле:

\begin{equation}
\label{eq:to3}
k_{\textsl{ч}} = \frac{I_{\textsl{k.мин.}}^{(2)}}{I_{\textsl{с.з.ТО}}},
\end{equation}
\begin{eqexpl}[15mm]
\item{$I_{\textsl{р.мин.}}$} минимальное значение тока при металлическом двухфазном КЗ на выводах ВН защищаемого трансформатора, А;
\item{$I_{\textsl{с.з.ТО}}$} ток срабатывания отсечки, вычисленный по условию (\ref{eq:to2}), А.
\end{eqexpl}

В соответствии с \cite{prikaz546} наименьший коэффициент чувствительности $k_{\textsl{ч}}$ для токовых защит трансформатора без выдержек времени должен быть равным 2,0.

\item
Если токовая отсечка не удовлетворяет требованиям чувствительности, а МТЗ имеет выдержку более 0,5 с, на трансформаторах мощностью более 1 МВА должна устанавливаться дифференциальная токовая защита.

\end{enumerate}